% Drawing 156-8/06-2 -- folder 156-8, batch 06, archive page 117 (PDF page 118).
% Transcription: en marge : un segment $\Phi$ formé de deux intervalles fermés raccordés en un point $(\bullet)$, « $\Phi$ non morcelée », et en dessous le même avec $\overline{\Phi}$, « morcelée ».
% Kind: segments joined at a point (marginal). Method 2 (schematic redraw from the facsimile). Status: ok.
% Notes: In the margin, written slantwise: Phi made of two closed intervals joined at a point (•), « Phi non morcelée », and below it the same with Phi-bar, « morcelée ». The slanted marginal words are given as two labels; the rest of the marginal note is text of the transcription.
% Coordinates: 1 unit = 1 mm = 1/100 of the crop width (bbox in page fractions: [0.08, 0.26, 0.27, 0.375]).
% Redrawn by Claude (claude-opus-5-5), issue #73 pilot, 2026-11-07; not checked by a human; not compiled (no TeX in the session).
% Rights: reproduces part of the work -- behind RENDER_TRANSCRIPTS, never published by itself.
\documentclass[tikz,border=3pt]{standalone}
\usepackage{amsmath,amssymb}
\usetikzlibrary{patterns,arrows.meta}
\tikzset{pencil/.style={black!45}, struck/.style={black!35}, mark/.style={circle,draw,inner sep=1.3pt}, pt/.style={circle,fill,inner sep=0.9pt}, opt/.style={circle,draw,fill=white,inner sep=1.1pt}}
\begin{document}
\begin{tikzpicture}[x=1mm,y=1mm,line cap=round,line join=round,font=\scriptsize]
\draw (20,-43.9) -- (38,-27.9);
\node[inner sep=1pt] at (39.5,-25.9) {\tiny $(\bullet)$};
\draw (42,-23.9) -- (55,-12.4);
\draw (18.5,-41.9) -- (21.5,-45.9);
\draw (53.5,-10.4) -- (56.5,-14.4);
\draw (27.5,-56.9) -- (43.5,-43.9);
\node[inner sep=1pt] at (47,-39.9) {\tiny $(\bullet)$};
\draw (50,-37.4) -- (62,-26.9);
\draw (26,-54.9) -- (29,-58.9);
\draw (60.5,-24.9) -- (63.5,-28.9);
\node[inner sep=1pt] at (29,-16.4) {\scriptsize $\Phi$};
\node[inner sep=1pt] at (52,-20.9) {\tiny non morcelée};
\node[inner sep=1pt] at (14.5,-68.9) {\scriptsize $\overline{\Phi}$};
\node[inner sep=1pt] at (30,-66.9) {\tiny morcelée};
\end{tikzpicture}
\end{document}
